% !TeX root = ../动手学强化学习笔记.tex
% 由主文档通过 \input 载入；请编译主文档。
\setcounter{chapter}{19} % 保留原书章号。
\chapter{部分可观测决策}
\label{ch:rl-pomdp}

部分可观测马尔可夫决策过程（partially observable Markov decision process，POMDP）中，环境状态仍可满足马尔可夫性质，智能体却不能直接获得它。合理行为可能依赖先前线索，因此策略需要利用历史或其信息摘要。本章承接条件概率与 MDP，重点推导有限隐状态、模型已知时的信念更新，再说明记忆策略与循环网络如何近似利用历史；学习目标是区分精确信息状态与学习得到的记忆表示。内容以建模、滤波和序列训练的基本要求为限，信念空间上的完整规划算法留待进一步学习。

\section{隐状态、观测核与历史}

本章采用联合转移奖励核 $p(s',r\mid s,a)$ 及观测核 $\Omega(o'\mid s')$，初始观测由相应初始状态生成。可用历史是 $H_t=(O_0,A_0,R_1,\ldots,R_t,O_t)$，策略为 $\pi(a\mid H_t)$。观测核还可推广为依赖动作的形式，但推导时必须保持统一约定。

两个隐状态可能产生同一观测而要求不同动作。例如迷宫的两个外观相同路口需要不同转向，单张观测无法区分，先前经过的标志却能确定位置。把观测直接当作 MDP 状态会丢失预测未来所需的信息；增大无记忆网络并不能恢复输入中根本不存在的信息。

\section{信念状态与贝叶斯滤波}

信念状态（belief state）为 $b_t(s)=\Pr(S_t=s\mid H_t)$。它把历史压缩成关于隐状态的后验分布。在有限状态且模型已知时，执行 $a$ 并观察奖励 $r$ 和新观测 $o'$ 后，
\begin{equation}
b_{t+1}(s')=
\frac{\Omega(o'\mid s')\sum_s b_t(s)p(s',r\mid s,a)}
{\sum_u\Omega(o'\mid u)\sum_s b_t(s)p(u,r\mid s,a)}.
\label{eq:rl-belief-update}
\end{equation}
分母是所观察证据的概率，要求为正；连续奖励使用相应密度。当奖励本身携带状态信息时不能将它丢弃。若奖励由已有信息确定而不提供额外证据，可退化为先用 $P$ 预测、再按观测似然归一化的标准形式。

\begin{example}[一次有噪观测]
隐状态为左或右，先验概率为 $(0.6,0.4)$。传感器报告“左”的概率在真实左状态下为 $0.8$、真实右状态下为 $0.2$，且本步状态不变、奖励不提供信息。观察“左”后，左状态后验为 $0.6\times0.8/(0.6\times0.8+0.4\times0.2)=6/7$。
\end{example}

给定 $b_t$ 和动作，后继证据分布与更新后的信念都由模型决定，因而信念是新的马尔可夫信息状态。即时期望奖励为 $\bar r(b,a)=\sum_s b(s)r(s,a)$，可以在信念空间建立贝尔曼方程。代价是即使原状态有限，信念也遍布概率单纯形，通常形成连续状态问题。

\section{记忆策略与循环网络}

显式滤波需要已知或学习的概率模型。另一种方法用记忆 $h_{t+1}=f_\theta(h_t,o_{t+1},a_t,r_{t+1})$，再由 $\pi_\theta(a\mid h_t)$ 决策。固定长度观测堆叠只保留最近窗口，循环网络原则上能保留更长历史，但训练到的 $h_t$ 不自动等于精确信念。

序列训练必须记录回合边界和初始隐藏状态。随机抽取中间片段时，可用较早观测预热记忆或保存一致的历史状态；若直接从零开始，训练时的记忆分布可能与真实执行不同。截断反向传播限制梯度传递长度，不意味着环境中的真实记忆依赖也被截断。

\section{习题与实验}
\begin{enumerate}
\item 在后验例子中连续观察两次“左”，假设状态不变且观测条件独立，求新后验；再解释加入状态转移后为何不能简单重复同一个静态公式。
\item 构造奖励能区分隐状态而观测不能的模型，比较使用和忽略奖励的信念更新。
\item 设计开始时显示一次线索、若干步后才要求选择的环境，比较无记忆、固定窗口和循环策略；用完整状态策略作为信息充足的参照。
\end{enumerate}





